% =====================================================================
%   附录 C：数据库 (Database)
% =====================================================================

\section{数据库}\label{app:db}

\subsection{数据集细节}

\subsubsection*{构造细节}

我们通过复用并整合多个已有数据集来获取源查询与数据库：WikiSQL
\citep{zhong2017}、WikiTableQuestions \citep{pasupat2015}、SQA
\citep{iyyer2017}、HybridQA \citep{chen2020hybridqa} 与 FeTaQA
\citep{nan2021fetaqa}，以确保指令与数据的多样性。

为进一步丰富数据集并避免潜在的数据泄露，我们采用 \code{gpt-3.5-turbo}
进行数据增强：给定表的表头信息与原始行，\code{gpt-3.5-turbo} 将生成 10 条
新的数据行；同时，利用表名、表头信息以及若干 SQL 示例，\code{gpt-3.5-turbo}
被要求额外生成 5 条 SQL 查询。每条得到的 SQL 语句被依次送入
\code{gpt-3.5-turbo}，要求其在不改变原意的前提下重新表述。各有效条目
经过筛选后，被采样进入最终的数据集（共 300 条），并被归为三类基本
数据库操作：\code{select}、\code{insert} 与 \code{update}。

因此，数据集中每条样本包含以下内容：
\begin{itemize}
    \item \textbf{指令（Instruction）}：用于描述问题并引导智能体动作的自然
          语言描述。
    \item \textbf{表信息（Table Info）}：关于表名与列名（即元信息）的说明。
    \item \textbf{表内容（Table Content）}：表中的实际内容，用于创建数据库。
    \item \textbf{正确答案（Correct Answer）}：对于 \code{select} 类样本，
          为文本答案；对于其它类型（即 \code{insert}、\code{update}）的
          样本，为被正确修改后表的哈希码。
\end{itemize}

\subsubsection*{评测设置}

我们按以下流程对数据集中的每条问题进行评估：
\begin{enumerate}
    \item \textbf{初始化}：根据表内容构造初始 SQL 脚本，并在 Docker 容器中
          初始化一个 MySQL 数据库，对外暴露一个转发端口用于交互。
    \item \textbf{交互}：初始提示引导智能体给出一条可执行的 SQL 命令及其
          推理过程。智能体接收提示、指令与表信息描述，并被期望以指定格式
          返回响应。我们执行该 SQL 语句并直接将结果返回给智能体，循环往复，
          直至智能体提交最终答案或发生错误（例如达到最大交互轮数或解析动作
          失败）。
    \item \textbf{校验}：对于 \code{select} 类问题，我们将智能体的答案与
          标准文本答案进行比较（忽略顺序，但要求完全匹配）。若答案为单一
          数字，则所有等价表示均被接受（例如 \code{5}、\code{"5.0"}、
          \code{'+5'} 视为相同）。对于 \code{insert} 或 \code{update} 类
          问题，我们计算并比较智能体执行操作后表的哈希值与正确 SQL 操作后
          表的哈希值。
\end{enumerate}

\subsubsection*{评测指标}

本文以智能体完成任务的成功率作为评测指标。总体成功率为三类任务成功率的
宏平均。

\subsection{数据增强}

我们对基于既有 SQL 数据集的三类 DB 任务的数据增强进行详细说明
\citep{zhong2017,pasupat2015,iyyer2017,chen2020hybridqa,nan2021fetaqa}，
这些数据集均为不包含 \code{insert}、\code{update} 等常见操作的 QA 类问题。
我们首先验证原始数据的有效性，然后从经过过滤的数据中按类随机采样，构成
最终数据集。我们采用 \code{gpt-3.5-turbo} 对原始指令进行扩充与改写。
\begin{itemize}
    \item \textbf{Insert}：给定表的表名、表头信息以及原始行，我们生成 5 条
          用于插入的 SQL 语句。随后在不改变原意的前提下对其进行改写
          （使用更短或更长的表达，或调整语序）。
    \item \textbf{Update}：给定表名、表头信息以及先前生成的 5 条插入 SQL
          语句，我们基于这些语句再生成 5 条用于修改的 SQL 语句，并按上述
          标准进行改写。
\end{itemize}

为保证数据质量，每条经增强的查询语句均需通过单元测试脚本。任务中的查询
类型均落在传统的 Text-to-SQL 评测范畴之内，我们仅对其进行采样与归类。
在既有数据集中的每条查询语句被分类为以下类型之一：\code{Counting}、
\code{Aggregation-MIN}、\code{Aggregation-MAX}、\code{Aggregation-AVG}、
\code{Aggregation-SUM}、\code{Ranking} 或 \code{Comparison}，每条只能归入
一种类型；其余则归入 \code{Other}。

\subsection{提示词示例}

我们使用如下格式的提示词：
\begin{quote}\small\ttfamily
User: I will ask you a question, then you should help me operate a MySQL
database with SQL to answer the question. You have to explain the problem
and your solution to me and write down your thoughts. After thinking and
explaining thoroughly, every round you can choose to operate or to answer.\\
your operation should be like this:\\
Action: Operation\\
\verb|```sql|\\
SELECT * FROM table WHERE condition;\\
\verb|```|\\
You MUST put SQL in markdown format without any other comments. Your SQL
should be in one line. Every time you can only execute one SQL statement.
I will only execute the statement in the first SQL code block. Every time
you write a SQL, I will execute it for you and give you the output.\\
If you are done operating, and you want to commit your final answer, then
write down:\\
Action: Answer\\
Final Answer: ["ANSWER1", "ANSWER2", \dots]\\
DO NOT write this pattern unless you are sure about your answer. I expect
an accurate and correct answer. Your answer should be accurate. Your answer
must be exactly the same as the correct answer. If the question is about
modifying the database, then after done operation, your answer field can
be anything. If your response cannot match any pattern I mentioned earlier,
you will be judged as FAIL immediately. Your input will be raw MySQL
response, you have to deal with it by yourself.
\end{quote}

\subsection{数据增强中的偏差研究}

为验证本文数据集未引入由模型在数据增强过程中所产生的偏差，我们使用
\code{claude-2} 重新标注了一小批量数据，并与以 \code{gpt-4} 和
\code{gpt-3.5-turbo} 为例的标注结果进行比较。

\begin{table}[h]
    \centering
    \small
    \caption{
        \textbf{表 5：在新标注数据上 \code{gpt-4} 与 \code{gpt-3.5-turbo}
        的得分对比（原标注数据 vs.\ 新标注数据）。}
    }
    \label{tab:db-bias}
    \begin{tabular}{l c c c c}
        \toprule
        \textbf{模型} & \textbf{数据来源} & \textbf{SELECT} & \textbf{INSERT} & \textbf{UPDATE} \\
        \midrule
        \code{gpt-4}          & 原标注 & 0.32 & 0.32 & 0.32 \\
        \code{gpt-3.5-turbo}  & 原标注 & 0.21 & 0.23 & 0.66 \\
        \code{gpt-4}          & 新标注 & ---  & 0.27 & 0.66 \\
        \code{gpt-3.5-turbo}  & 新标注 & ---  & 0.19 & 0.92 \\
        \bottomrule
    \end{tabular}
\end{table}

如表~\ref{tab:db-bias} 所示，数据在两类标注下表现出相似的得分模式——
\code{gpt-4} 在 UPDATE 类操作上表现较弱，而在 INSERT 类任务上则更胜一筹。
这一观察表明，本文的数据增强方法并未引入明显的偏差，不同操作间的相对
得分关系得以保持。
